\scanpage{6}
\mspage{P 52}
cette ``promenade'' (dans l'étape ``Points de vue et vision'', n$^\circ$ 6), ces théorèmes ne prennent pour moi tout leur sens que par le contexte nourricier d'un grand thème, initié par une de ces ``idées fécondes''. Leur démonstration dès lors découle, comme de source et sans effort, de la nature même, de la ``profondeur'' du thème qui les porte - comme les vagues du fleuve semblent naître en douceur de la profondeur même de ses eaux, sans rupture et sans effort. Je m'exprime dans un sens tout analogue, mais avec d'autres images, dans la note déjà citée ``La mer qui monte...''.

\noindent\rule{2.1cm}{0.4pt}

18.\quad Quand cet ``avant-propos'' a commencé à tourner à la promenade à travers mon oeuvre de mathématicien, avec mon petit topo sur les ``héritiers'' (bon teint) et sur les ``bâtisseurs'' (incorrigibles), a commencé aussi à apparaître un n o m pour cet avant-propos manqué : ce serait ``L'enfant et le bâtisseur''. Au cours des jours suivants, il devenait de plus en plus clair que ``l'enfant'' et ``le bâtisseur'' étaient un seul et même personnage. Ce nom est donc devenu, plus simplement, ``L'enfant bâtisseur''. Un nom, ma foi, qui ne manquait pas d'allure, et tout fait pour me plaire !

Mais voilà que la réflexion fait apparaître que cet altier ``bâtisseur'', ou (plus modestement) l'enfant-qui-joue-à-faire-des-maisons, ce n'était qu'un des visages du fameux enfant-qui-joue, lequel en avait d e u x. Il y a aussi l'enfant-qui-aime-à-explorer-les-choses, à aller fouiner et s'enfouir dans les sables ou dans les vases boueuses et sans nom, les endroits les plus impossibles et les plus saugrenus... Pour donner le change sans doute (ne serait-ce qu'à moi-même...), j'ai commencé par l'introduire sous le nom flamboyant de ``pionnier'', suivi de celui, plus terre-à-terre mais encore auréolé de prestige, d' ``explorateur''. C'était à se demander, entre le ``bâtisseur'' et le ``pionnier-explorateur'', lequel était le plus mâle, le plus alléchant des deux ! Pile ou face ?

Et puis, en y regardant d'un peu plus près, voilà notre intrépide ``pionnier'' qui se trouve finalement être une f i l l e (qu'il m'avait plu d'habiller en garçon) - une soeur des mares, de la pluie, des brumes et de la nuit, silencieuse et quasiment invisible à force de s'effacer dans l'ombre - celle que toujours on oublie (quand on ne fait mine de se gausser d'elle...). Et j'ai bien trouvé moyen moi aussi, pendant des jours et des jours, de l'oublier - de l'oublier doublement, pourrais-je dire : je n'avais voulu voir d'abord que le garçon (celui qui joue à faire des maisons...) - et même quand je n'ai pu m'empêcher, à force, de voir quand même l ' a u t r e , je l'ai vue encore en garçon, elle aussi...

Pour ce qui est du beau nom pour ma promenade, du coup il ne tient plus du tout. C'est un nom tout-en-yang, tout ``macho'', un nom-qui-boite. Pour le faire tenir pas de guingois, il faudrait y faire figurer l ' a u t r e également. Mais, chose étrange, `` l ' a u t r e ''\quad n ' a \quad p a s \quad v r a i m e n t \quad d e
% archive sheet number: 5

